\subsection{Requerimientos del ambiente}

Tras haber finalizado el proceso de migración desde CUDA a los estándares que serán analizados, se debe establecer los requerimientos que debe cumplir este ambiente para asegurar la correcta ejecución y condiciones lo más idénticas posibles entre entornos. Para conseguir esto, partimos del uso de Docker, debido a su funcionamiento similar a aquel de las máquinas virtuales, donde se asegura la ejecución del código en un ambiente hermético sin el sobrecosto de recursos y configuración que requiere una máquina virtual.

Para evitar conflictos entre configuraciones, se usará una imágen de Docker diferente por cada entorno, cada uno de ellos tiene un proceso de configuración inicial e instalación diferente, la separación entre entornos asegura la minimización de inconvenientes que puedan ocurrir durante estos procesos. Ahora bien, Docker nos permite crear una imágen con base en una imágen previa, con esta funcionalidad, se usará como imágen base, la imágen correspondiente al sistema operativo Ubuntu, en su última versión LTS (Long Term Support), que es su versión más estable y actualizada. Esta decisión fue tomada en primera instancia, debido a la familiaridad y popularidad de esta distribución de Linux, su uso durante los procesos de migración, depuración, compilación y ejecución de las bases de código en todos los estándares en la etapa anterior de este estudio. Y en segunda instancia a la compatibilidad que tiene Ubuntu en la instalación y configuración requerida para estos ambientes,requiriendo la menor cantidad de intervención.

En adición a esto, para seguir el estándar de separación, cada imágen debe realizar 3 acciones: instalación y configuración inicial del estándar, que abarca la instalación de todas las dependencias, establecimiento de variables de entorno e instalación de los paquetes principales; compilación y ejecución de la base de código, donde se usará el hardware provisto para cada imágen para asegurar su correcta ejecución; recopilación de datos obtenidos, se tomarán todas las métricas que son de interés para el estudio y se almacenarán en un archivo de salida.

Dentro de las últimas 2 acciones, interviene un script cuyo trabajo es compactarlas en un archivo bash ejecutable, exclusivo de cada imágen. En este script, se ejecutará el código modificando las variables de entrada \textbf{num-seeds} y \textbf{num-steps}, 20 veces, manteniendo una variable constante mientras la otra aumenta en incrementos del 10\%, 10 iteraciones por variable, donde \textbf{num-steps} tiene un valor base de 1000 y \textbf{num-seeds} de 10.000. Esta ronda de pruebas se repetirá 10 veces, para un total de 200 ejecuciones. En cada una de estas ejecuciones se tomarán las medidas de tiempo de ejecución, medido internamente en el código, los valores de las dos variables en la iteración y los valores del reporte de \textit{valgrind} para el análisis del uso de memoria.



